\section{Source and Preprocessing Provenance}

\paragraph{Open3DSG.}
The main source-output case study uses the official non-averaged BLIP checkpoint
\texttt{epoch=13-step=13104.ckpt}. It was selected by train/dev validation loss,
not by the held-out GeoCalib result. The documented recovery route covers all
548 evaluation contexts by setting the minimum visible-object count to two and
regenerating relaxed views for two scans. The unmodified public-source route
covers 533/548 contexts. An older 127-scan workflow covered 377/388 contexts;
its corrected comparison covers 388/388 and changes Recall@100 by 0.28
percentage points. These older branches are provenance/sensitivity evidence,
not separate main results or candidate selection criteria.

\paragraph{SGFN.}
We use the official \texttt{full\_l160} source with its 160-object,
26-relation checkpoint. A pre-execution audit corrected the dataset-list name
and an incompatible \texttt{l20} checkpoint URL. The evaluated scans,
denominator, methods, and tensor-shape checks are recorded in the submitted
artifacts.
